% Calculus I review game, converted from the original slides.
% The \response is the answer from the slides; \solution adds a short
% worked solution.
\documentclass{article}
% Which PDF to make (the web game ignores this):
%   questions  clues only, e.g. a study sheet for students
%   answers    clues and responses
%   solutions  clues, responses, and worked solutions
% Add any of these, separated by commas, e.g. [solutions, writesay]:
%   writesay   show each clue's [write]/[say] mark (for you, not students)
%   board      always show the board overview (even with questions)
%   noboard    never show the board overview
\usepackage[solutions,writesay]{../calculatedrisk}

\gametitle{Sample Single Variable Calculus Game}

\begin{document}

%======================================================================
\begin{category}{Derivatives}

  \begin{clue}[write]
    It is the derivative of
    \[ f(x) = 4x^3 - 2x^2 + 9x - 7. \]
  \response
    What is $f'(x) = 12x^2 - 4x + 9$?
  \end{clue}

  \begin{clue}[write]
    It is the derivative of
    \[ g(x) = \frac{2}{x} - \frac{5x}{x^4}. \]
  \response
    What is $g'(x) = -\dfrac{2}{x^2} + \dfrac{15}{x^4}$?
  \end{clue}

  \begin{clue}[write]
    It is the derivative of
    \[ y = 4x^{5/2}. \]
  \response
    What is $\dfrac{dy}{dx} = 10x^{3/2}$?
  \end{clue}

  \begin{clue}[write]
    It is the derivative of
    \[ f(x) = \sqrt[3]{x} + 2x^7 - 2x^{-1/5}. \]
  \response
    What is $f'(x) = \dfrac{1}{3}x^{-2/3} + 14x^6 + \dfrac{2}{5}x^{-6/5}$?
  \end{clue}

  \begin{clue}[say]
    It is the derivative of
    \[ f(x) = \frac{x}{|x|}. \]
  \response
    What is $0,$ except at $x = 0,$ where it is undefined?
  \solution
    \[
    f(x) = \begin{cases}
      1, & x > 0, \\
      -1, & x < 0, \\
      \text{undefined}, & x = 0
    \end{cases}
    \]
    and
    \[
    f'(x) = \begin{cases}
      0, & x > 0, \\
      0, & x < 0, \\
      \text{undefined}, & x = 0
    \end{cases}
    \]
    $f$ and $f'$ are undefined at $x = 0.$
  \end{clue}

\end{category}

%======================================================================
\begin{category}{Limits}

  \begin{clue}[say]
    It is the value of
    \[ \lim_{x\to 3} \frac{x-3}{x^2-9}. \]
  \response
    What is $\dfrac{1}{6}$?
  \solution
    \[ \lim_{x\to 3} \frac{x-3}{(x-3)(x+3)} = \lim_{x\to 3} \frac{1}{x+3} = \frac{1}{6}. \]
  \end{clue}

  \begin{clue}[say]
    It is the value of
    \[ \lim_{x\to 3^-} \frac{x}{x-3}. \]
  \response
    What is $-\infty$?
  \end{clue}

  \begin{clue}[say]
    It is the value of
    \[ \lim_{x\to 0} \frac{\sin x}{x}. \]
  \response
    What is $1$?
  \end{clue}

  \begin{clue}[say]
    It is the value of
    \[ \lim_{x\to 1} \frac{x^2-1}{x^2-4x+3}. \]
  \response
    What is $-1$?
  \solution
    \[ \lim_{x\to 1} \frac{(x-1)(x+1)}{(x-1)(x-3)} = \lim_{x\to 1} \frac{x+1}{x-3} = \frac{2}{-2} = -1. \]
  \end{clue}

  \begin{clue}[say]
    It is the value of
    \[ \lim_{h\to 0} \frac{\sin(\pi+h) - \sin \pi}{h}. \]
  \response
    What is $-1$?
  \solution
    This is the definition of the derivative of $\sin x$ at $x = \pi,$ so the
    limit is $\cos \pi = -1.$
  \end{clue}

\end{category}

%======================================================================
\begin{category}{Definitions and Theorems}

  \begin{clue}[write]
    It is the definition of continuity.
  \response
    What is: $f$ is continuous at $x = a$ provided
    \[ \lim_{x\to a} f(x) = f(a)? \]
  \end{clue}

  \begin{clue}[say]
    It is the theorem we would use to show that $\cos x = x$ for some $x.$
  \response
    What is the Intermediate Value Theorem?
  \solution
    Let $g(x) = \cos x - x,$ which is continuous. Since $g(0) = 1 > 0$ and
    $g(\pi/2) = -\pi/2 < 0,$ the Intermediate Value Theorem gives a $c$ in
    $(0, \pi/2)$ with $g(c) = 0,$ that is, $\cos c = c.$
  \end{clue}

  \begin{clue}[write]
    It is the definition of the derivative.
  \response
    What is
    \[ f'(x) = \lim_{h\to 0} \frac{f(x+h) - f(x)}{h}? \]
  \end{clue}

  \begin{clue}[say]
    It is the theorem we would use to prove that
    \[ \lim_{x\to 0} x^2 \sin\left(\frac{1}{x}\right) = 0. \]
  \response
    What is the Squeeze (or Sandwich) Theorem?
  \solution
    Since $-1 \le \sin(1/x) \le 1,$ we have
    $-x^2 \le x^2 \sin(1/x) \le x^2.$ Both $-x^2$ and $x^2$ approach $0$ as
    $x \to 0,$ so the limit is $0.$
  \end{clue}

  \begin{clue}[write]
    It is the definition of the limit.
  \response
    What is: $\displaystyle\lim_{x\to a} f(x) = L$ if for every
    $\varepsilon > 0$ there exists $\delta > 0$ such that
    $|f(x) - L| < \varepsilon$ whenever $0 < |x - a| < \delta$?
  \end{clue}

\end{category}

%======================================================================
\begin{category}{More Derivatives}

  \begin{clue}[write]
    It is the derivative of $\csc x.$
  \response
    What is $-\csc x \cot x$?
  \end{clue}

  \begin{clue}[write]
    It is the derivative of $x^2 \cos x.$
  \response
    What is $2x\cos x - x^2 \sin x$?
  \end{clue}

  \begin{clue}[write]
    It is the derivative of
    \[ h(t) = \frac{3t+2}{2t+3}. \]
  \response
    What is $\displaystyle h'(t) = \frac{3(2t+3) - 2(3t+2)}{(2t+3)^2} = \frac{5}{(2t+3)^2}$?
  \end{clue}

  \begin{clue}[doubledown, write]
    It is the derivative of
    \[ f(x) = \cot(x^2 + x). \]
  \response
    What is $\dfrac{df}{dx} = -(2x+1)\csc^2(x^2+x)$?
  \end{clue}

  \begin{clue}[write]
    It is the derivative of $\tan(\tan(\tan x)).$
  \response
    What is $\sec^2(\tan(\tan x)) \sec^2(\tan x) \sec^2 x$?
  \end{clue}

\end{category}

%======================================================================
\begin{category}{Applications of the Derivative}

  \begin{clue}[write]
    It is the tangent line to the curve
    $y = x^3 - \sqrt{x} $ at $x = 1.$
  \response
    What is $y = \frac{5}{2}(x - 1)$?
  \solution
    At $x = 1,$ $y = 1 - 1 = 0.$ The slope is
    \[ y'(1) = 3(1)^2 - \frac{1}{2\sqrt{1}} = \frac{5}{2}. \]
  \end{clue}

  \begin{clue}[say]
    They are the locations of the horizontal tangent lines to the graph of
    $ y = 2x^3 - 12x^2.$
  \response
    What are $x = 0$ and $x = 4$?
  \solution
    \[ y' = 6x^2 - 24x = 6x(x - 4) = 0 \]
    when $x = 0$ or $x = 4.$
  \end{clue}

  \begin{clue}[say]
    They are the units of $dm/dt,$ where $m(t)$ is the number of m\&m's
    eaten in $t$ minutes.
  \response
    What are m\&m/m?
  \end{clue}

  \begin{clue}[say]
    It is the initial velocity, $v(0),$ of the particle with position function
    \[ p(t) = 3\cos(t + \pi). \]
  \response
    What is $0$?
  \solution
    \[ v(t) = p'(t) = -3\sin(t + \pi), \]
    so $v(0) = -3\sin \pi = 0.$
  \end{clue}

  \begin{clue}[say]
    It is the tangent line to $y = \sqrt[3]{x}$ at $x = 0.$
  \response
    What is the line $x = 0$?
  \solution
    $y' = \frac{1}{3}x^{-2/3},$ which goes to $\infty$ as $x \to 0.$ The
    tangent line at the origin is vertical: the $y$-axis, $x = 0.$
  \end{clue}

\end{category}

%======================================================================
\begin{category}{Potpourri}

  \begin{clue}[say]
    It is an example of a continuous function that is not differentiable.
  \response
    There are many responses. For example, what is $f(x) = |x|$?
  \end{clue}

  \begin{clue}[say]
    He is one of the inventors of calculus.
  \response
    Who is Isaac Newton or Gottfried Leibniz?
  \end{clue}

  \begin{clue}[say]
    He is the other inventor of calculus.
  \response
    Who is Gottfried Leibniz or Isaac Newton, whichever wasn't named
    before?
  \end{clue}

  \begin{clue}[say]
    It is the derivative of $\pi^e + e^\pi.$
  \response
    What is $0$?
  \end{clue}

  \begin{clue}[say]
    It is Cady Heron's response to the final question of the math
    competition in \emph{Mean Girls}, after she sees that the one-sided
    limits disagree.
  \response
    What is ``the limit does not exist''?
  \end{clue}

\end{category}

%======================================================================
\begin{final}[write]{Pathological Examples}
  True or false: $f'(0)$ exists for the function
  \[
  f(x) =
  \begin{cases}
    x^2 & \text{if } x \text{ is rational,} \\
    0 & \text{if } x \text{ is irrational.}
  \end{cases}
  \]
\response
  What is true?
\solution
  Since $f(0) = 0,$
  \[ f'(0) = \lim_{h\to 0} \frac{f(h)}{h}, \]
  and $\left|\dfrac{f(h)}{h}\right| \le \dfrac{h^2}{|h|} = |h| \to 0.$ By the
  Squeeze Theorem, $f'(0) = 0.$
\end{final}

\end{document}
